%% intro -- the problem, its history, the main theorem and the bounds, the
%% mechanization, an overview of the uniform proof route with pointers to the
%% sections, the running examples, contributions, and the map of the paper.

\providecommand{\parity}{\operatorname{par}}
\providecommand{\CondI}{\textup{(I0)}}
\providecommand{\CondU}{\textup{(U)}}
\providecommand{\CondII}{\textup{(II)}}

\section{Introduction}\label{sec:intro}

Type inference with subtyping expresses the requirements of a program as
subtype constraints between type expressions with unknowns
\citep{Mitchell91,FuhMishra90,AikenWimmers93,EifrigSmithTrifonov95,Pottier98}.
An inferred \emph{constrained type} $\forall\bar\alpha.\,C\Rightarrow\tau$
pairs a type $\tau$ with a conjunction $C$ of constraints; its instances are
the types $\rho(\tau)$ for the solutions $\rho$ of $C$
\citep{TrifonovSmith96}. Inferred constraint sets grow quickly, and practical
systems simplify them
\citep{Pottier96,Pottier01,TrifonovSmith96,FahndrichAiken96,AikenWimmersPalsberg99,FlanaganFelleisen99}.
Whether a simplification is correct, and whether one constrained type may be
used where another is expected, are questions about the solution sets of
constraint systems.

The simplest such question is \emph{entailment}: does every solution of a
conjunction $\varphi$ satisfy an inequality $x\le y$? Constrained types,
however, also involve variables that occur on one side only. Let
$\varphi(\bar x)$ and $\psi(\bar x,\bar z)$ be finite conjunctions of subtype
constraints, where the \emph{hidden} variables $\bar z$ do not occur in
$\varphi$. \emph{Existential entailment} $\varphi\entails\Exists\psi$ holds if
every solution of $\varphi$ extends to a solution of $\psi$ by some choice of
trees for $\bar z$ (\cref{def:bg-existential}); logically, it is the validity
of $\forall\bar x\,(\varphi\Rightarrow\Exists\psi)$. Plain entailment is the
case without hidden variables, and satisfiability of $\psi$ is the case in
which $\varphi$ is empty and every variable is hidden.

\paragraph{Subsumption and Simplification}
\citet{TrifonovSmith96} define subtyping between constrained types
semantically: $\forall\bar\alpha.\,C_1\Rightarrow\tau_1$ is below
$\forall\bar\beta.\,C_2\Rightarrow\tau_2$ if every instance of the right-hand
type is a supertype of some instance of the left-hand type. With $\bar\alpha$
and $\bar\beta$ renamed apart from each other and from the free variables of
both types (\cref{prop:ap-subsumption}), this is the existential entailment
\begin{equation}\label{eq:intro-subsumption}
  \forall\bar\alpha.\,C_1\Rightarrow\tau_1
  \;\le\;
  \forall\bar\beta.\,C_2\Rightarrow\tau_2
  \quad\iff\quad
  C_2\entails\exists\bar\alpha\,(C_1\wedge\tau_1\le\tau_2),
\end{equation}
in which the variables of the right-hand type are universally quantified and
the witnesses belong to the smaller, more general type. For example,
$\forall\alpha.\,\alpha\to\alpha\le\beta\to\beta$ holds with the witness
$\alpha=\beta$, while the converse fails; \cref{prop:ap-subsumption} adds the
assumptions on program variables of \citet{TrifonovSmith96}. Likewise,
replacing a system $C(\bar x,\bar z)$ by $D(\bar x,\bar k)$, where $\bar x$
remain observable, is correct when $C\entails\exists\bar k\,D$ and
$D\entails\exists\bar z\,C$. Removing a redundant constraint needs only plain
entailment; eliminating or merging internal variables needs the existential
form (\cref{sec:ap}). The arrow is contravariant in its domain, so a variable
of a constrained type of a higher-order program typically occurs both co- and
contravariantly, as $\alpha$ does in $\alpha\to\alpha$.

\paragraph{Why Plain Entailment Does Not Suffice}
Consider trees built from a bottom element~$\Bot$ and a binary covariant
constructor~$f$, and the consequent
\[
  \psi \;=\; x\le z \,\wedge\, z\le y \,\wedge\, z\le f(z,z).
\]
The literal $z\le f(z,z)$ says that all subtrees of $z$ lie below $z$. Write
$x\cdot\eta$ for the subtree of $x$ at the path $\eta\in\{0,1\}^*$, with
$x\cdot\eta=\Bot$ when $x$ has no node at $\eta$. Then $\exists z\,\psi$ holds
exactly when $x\cdot\eta\le y$ for every path $\eta$: necessity follows from
$x\cdot\eta\le z\cdot\eta\le z\le y$, and the least upper bound of all the
subtrees $x\cdot\eta$ is a witness. Thus $\varphi\entails\exists z\,\psi$ is an
infinite family of plain entailments, and the witness for $z$ depends on the
solution of $\varphi$. Sound tests replace hidden variables by canonical
bounds or search for syntactic derivations
\citep{TrifonovSmith96,Pottier01}; they are incomplete, and \cref{sec:ap}
recalls a correct simplification of \citet{Pottier98thesis} that such a test
misses.

\begin{example}[Shift closure: the running example]\label{ex:intro-shift}
Take $\psi$ as above, with $x,y$ visible and $z$ hidden. Let
$\Kpos(x)\subseteq\{0,1\}^*$ be the set of positions of the $f$-nodes of $x$,
so that $x\le y$ iff $\Kpos(x)\subseteq\Kpos(y)$ (\cref{lem:bg-bf-order}), and
let $\Suf(L)$ be the set of suffixes of words of~$L$. Since
$\Kpos(x\cdot\eta)=\eta^{-1}\Kpos(x)$, the condition ``$x\cdot\eta\le y$ for
every $\eta$'' says that $\exists z\,\psi$ holds iff
$\Suf(\Kpos(x))\subseteq\Kpos(y)$. For $\varphi_1=x\le y$ the entailment fails:
the solution $x=y=f(f(\Bot,f(\Bot,\Bot)),\Bot)$ of $\varphi_1$, with
$\Kpos(x)=\{\eps,0,01\}$, has the position $01$, whose suffix $1$ is not in
$\Kpos(y)$. For $\varphi_2=x\le y\wedge x\le f(x,b)\wedge b=\Bot$, the
entailment holds: $z=x$ is a witness, because $x\le f(x,b)\le f(x,x)$. For
$\varphi_3=x\le y\wedge x\le f(f(\Bot,f(\Bot,\Bot)),\Bot)$ the entailment
fails with the same countermodel as for $\varphi_1$, in which $x$ takes the
largest value that $\varphi_3$ allows.
\end{example}

\Cref{fig:intro-shift} draws two solutions of $\varphi_1$. For the first, the
join of all subtrees of $x$ fits below $y$ and is the least witness for $z$;
for the second it does not, and no witness exists. The choice of $z$ is made
after $x$ and $y$ are known, from their values.

\begin{figure}[tbp]
\centering
\begin{tikzpicture}[>=Stealth,font=\small,
    every node/.style={inner sep=1.5pt},
    var/.style={draw,circle,minimum size=6.5mm,inner sep=0pt},
    hid/.style={var,dashed},
    tree/.style={level distance=6.5mm,
      level 1/.style={sibling distance=11mm},
      level 2/.style={sibling distance=6mm},
      level 3/.style={sibling distance=5mm}},
    hl/.style={draw,rounded corners=1pt,fill=black!12},
    lb/.style={font=\footnotesize,align=center},
    e0/.style={left,font=\scriptsize,inner sep=1pt,pos=0.45},
    e1/.style={right,font=\scriptsize,inner sep=1pt,pos=0.45}]
  % the consequent: visible x, y and hidden z
  \node[var] (x) at (0,0.2) {$x$};
  \node[hid] (z) at (1.35,0.2) {$z$};
  \node[var] (y) at (2.7,0.2) {$y$};
  \draw[->] (x) -- node[above]{$\le$} (z);
  \draw[->] (z) -- node[above]{$\le$} (y);
  \path (z) edge[->,loop below,min distance=7mm,in=-120,out=-60] (z);
  \node[lb] at (1.35,-0.85) {$z\le f(z,z)$};
  \node[lb,text width=3.3cm] at (1.35,-1.75)
    {$x,y$ visible: fixed by\\ a solution of $\varphi$;\\
     $z$ hidden: chosen after};
  % row (a)
  \node[anchor=east,font=\small] at (4.1,1.9) {(a)};
  \begin{scope}[tree]
    \node (a1) at (5.4,2.5) {$f$}
      child {node {$f$} child {node {$\Bot$} edge from parent node[e0]{0}}
                        child {node {$\Bot$} edge from parent node[e1]{1}}
             edge from parent node[e0]{0}}
      child {node {$\Bot$} edge from parent node[e1]{1}};
    \node (a2) at (10.2,2.5) {$f$}
      child {node {$f$} child {node {$\Bot$} edge from parent node[e0]{0}}
                        child {node {$\Bot$} edge from parent node[e1]{1}}
             edge from parent node[e0]{0}}
      child {node {$\Bot$} edge from parent node[e1]{1}};
  \end{scope}
  \node[lb,anchor=south] at (5.4,2.65) {$x=y$};
  \node[lb,anchor=south] at (10.2,2.65) {least witness $z=x$};
  \node[lb] at (5.4,0.8) {$\Kpos(x)=\{\eps,0\}$};
  \node[lb] at (10.2,0.8) {$\Kpos(z)=\{\eps,0\}$};
  \draw[->] (6.6,1.9) -- node[above,lb]{join of all} node[below,lb]{subtrees $x\cdot\eta$} (9.0,1.9);
  % row (b)
  \node[anchor=east,font=\small] at (4.1,-1.1) {(b)};
  \begin{scope}[tree]
    \node (b1) at (5.4,-0.5) {$f$}
      child {node {$f$} child {node {$\Bot$} edge from parent node[e0]{0}}
                        child {node {$f$} child {node {$\Bot$} edge from parent node[e0]{0}}
                                          child {node {$\Bot$} edge from parent node[e1]{1}}
                               edge from parent node[e1]{1}}
             edge from parent node[e0]{0}}
      child {node[hl] (b1r) {$\Bot$} edge from parent node[e1]{1}};
    \node (b2) at (10.2,-0.5) {$f$}
      child {node {$f$} child {node {$\Bot$} edge from parent node[e0]{0}}
                        child {node {$f$} child {node {$\Bot$} edge from parent node[e0]{0}}
                                          child {node {$\Bot$} edge from parent node[e1]{1}}
                               edge from parent node[e1]{1}}
             edge from parent node[e0]{0}}
      child {node[hl] (b2r) {$f$} child {node {$\Bot$} edge from parent node[e0]{0}}
                                  child {node {$\Bot$} edge from parent node[e1]{1}}
             edge from parent node[e1]{1}};
  \end{scope}
  \node[lb,anchor=south] at (5.4,-0.35) {$x=y$};
  \node[lb,anchor=south] at (10.2,-0.35) {least candidate for $z$};
  \node[lb] at (5.4,-2.9) {$\Kpos(x)=\{\eps,0,01\}$};
  \node[lb] at (10.2,-2.9) {$\Kpos(z)=\{\eps,0,1,01\}$};
  \draw[->] (6.6,-1.1) -- node[above,lb]{join of all} node[below,lb]{subtrees $x\cdot\eta$} (9.0,-1.1);
  \node[lb,anchor=west] at (11.5,-1.25) {$\not\le y$};
\end{tikzpicture}
\caption{The shift closure of \cref{ex:intro-shift} over $\BF$. A solution of
$\varphi$ fixes the visible variables $x$ and $y$; the hidden $z$ is chosen
afterwards and depends on it. The literal $z\le f(z,z)$ makes the
constructor positions of $z$ closed under suffixes, so $z$ must lie above every
subtree of $x$, and their join is the least candidate. (a) For
$x=y=f(f(\Bot,\Bot),\Bot)$ the candidate is $x$ itself, which fits below $y$.
(b) For $x=y=f(f(\Bot,f(\Bot,\Bot)),\Bot)$, the countermodel of $\varphi_1$, the
position $01$ of $x$ forces a constructor of $z$ at its suffix $1$ (shaded),
where $y$ has the leaf $\Bot$; no $z$ exists.}
\Description{On the left, three circles x, z and y joined by arrows labeled
less or equal, with a loop at z for z below f of z and z; the circle of z is
dashed because z is hidden. On the right, two rows. In row a, the tree
f of f of bottom and bottom, and bottom, which is both x and y, is joined by an arrow
labeled join of all subtrees to the same tree, the least witness for z. In
row b, the tree of x and y has constructor positions epsilon, 0 and 01, and
the arrow leads to a tree with an extra constructor at position 1, shaded
together with the leaf bottom of y at position 1; the least candidate for z is not below
y.}
\label{fig:intro-shift}
\end{figure}


\paragraph{An Open Problem}
In non-structural subtyping, $\Bot\le t\le\Top$ relates trees of different
shapes. Plain non-structural subtype entailment was posed as an open problem
in 1996 \citep{Rehof96TLCA}, and \citet{Zhao26} decides it for one
constructor of any arity and variance. The existential
problem has been open since the same year. \citet[Section~1]{TrifonovSmith96}
``leave open the question whether it is decidable'' and decide a sound
approximation. \citet[Section~3.2]{FahndrichAiken96} write that ``no complete
algorithm is known''. \citet[Section~7.2.2]{Su02} writes that ``for
existential entailment and subtyping constrained types, virtually nothing is
known at this point, except the lower bounds carried over from entailment'',
and that for non-structural subtyping the problem ``may in fact be
undecidable'', because an alternative undecidability proof for the
first-order theory appears to use a single quantifier alternation.
\citet[Section~1.4.2]{Priesnitz04} calls subtyping of constrained types a
``prominent open problem''. The first-order theory is indeed undecidable,
already over $\BF$ with $f$ binary and covariant, in all three tree classes
\citep{SuAikenNiehrenPriesnitzTreinen02}, but with six quantifier blocks and a
matrix with disjunctions and negations, whereas existential entailment has one
alternation, and its antecedent and consequent are conjunctions. Over feature
trees, existential entailment is decided by an inclusion of finite automata
\citep{MullerNiehrenTreinen98,MullerNiehrenTreinen01}, in a setting where a
fresh feature is always available, which a constructor of fixed arity does not
offer (\cref{sec:rel}).

\paragraph{Our Result}
We fix one constructor $f$ of arity $n\ge0$ and a \emph{variance}
$c:\{0,\ldots,n-1\}\to\{0,1\}$, which makes each argument covariant or
contravariant; the arrow is $n=2$ with $c(0)=1$ and $c(1)=0$
(\cref{sec:bg-trees}). Trees may be infinite unless stated otherwise.

\begin{theorem}[Main theorem]\label{thm:intro-main}
For every arity $n\ge0$ and every variance $c$, existential entailment
$\varphi\entails\Exists\psi$ over $\BFT$ is decidable over arbitrary, over
regular, and over finite trees. Over $\BF$, the same holds when $f$ is
covariant. If $\varphi$ and $\psi$ are satisfiable in the tree class considered
and the entailment fails, then the failure is certified by a
\emph{counterexample word} of length at most $E_3(O(s))$: a word over the
argument positions of $f$, or two such words joined by a separator, from which
a countermodel is computed (\cref{def:as-certificate}). Here
$E_3(x)=2^{2^{2^x}}$, $s$ is the number of variables after flattening, every
introduced name and every hidden variable included, and the constant does not
depend on $n$.
\end{theorem}

\Cref{sec:as} defines counterexample words (\cref{def:as-certificate}),
states the theorem precisely (\cref{thm:as-main}), shows that every
counterexample word yields a countermodel effectively
(\cref{lem:as-countermodel}), and assembles the proof. The theorem needs no restriction on the polarity of
the variables: a variable may occur both co- and contravariantly, in the
antecedent and in the consequent. In particular, it decides the subsumption of
\citet{TrifonovSmith96} between constrained types over $\Bot$, $\to$ and
$\Top$, and the correctness of simplifications of such types
(\cref{sec:ap}). We do not claim the signature $\BF$ with a contravariant
argument; \cref{rem:tc-gadget-needed} explains why the reduction to $\BFT$
needs covariance.

\paragraph{Bounds and Complexity}
\Cref{prop:as-bounds} and \cref{sec:cx} refine the theorem. Only two kinds of
failure, both described in \cref{sec:intro-overview}, have triply exponential
bounds: failures of condition \CondU\ and \emph{mortal holes}. The larger
bound belongs to the mortal holes, the only case that needs a new compression
technique; all other failures have singly exponential witnesses. The decision procedure does not enumerate words: a
nondeterministic search builds a witness level by level, keeping only a
summary of each level, and runs in 2-EXPSPACE, hence in 3-EXPTIME
(\cref{prop:cx-search,thm:cx-upper}); for plain entailment this saves one
exponential in space over the procedure of \citet{Zhao26} (\cref{rem:cx-plain}). For a unary
constructor, $n=1$, witnesses are singly exponential and the problem is in
PSPACE, for both variances and all three tree classes (\cref{thm:cx-unary}),
and for $n=0$ it is a finite enumeration (\cref{rem:cx-nullary}). On the
other side, shortest counterexample words can have length $2^{\Omega(s)}$
(\cref{prop:cx-length}), and the problem is PSPACE-hard for every fixed arity
$n\ge2$, already with $f$ covariant, by explicit reductions that use the additional arguments of $f$; bare
padding of a binary construction does not suffice (\cref{thm:cx-hard} and
\cref{ex:cx-padding}). The gap between PSPACE and 2-EXPSPACE, the gap between
singly and triply exponential witnesses, and hardness for $n=1$ remain open.

\paragraph{Mechanization}
The decidability results are checked in the Lean~4 proof assistant
\citep{MouraUllrich21} with Mathlib \citep{Mathlib20}: over $\BFT$ for every
arity, arity zero included, every variance and all three tree classes, and
over $\BF$ with $f$ covariant (\cref{sec:mechanization}). Each of the six
procedures is a computable definition with a correctness theorem without
premises, and plain entailment is decided as the special case without hidden
variables. The anchored compression of \cref{sec:mo} is mechanized with an
elementary bound of the same form as ours, and all search bounds are explicit
functions of the numbers of automaton states. The procedures take flat
literals as input. The flattening of arbitrary type expressions
(\cref{lem:bg-flattening}), the bounds in terms of the number of variables, the
construction of countermodels and the complexity results are not mechanized.

\subsection{Overview of the Proof}\label{sec:intro-overview}

The proof follows a single route for all arities and variances
(\cref{fig:intro-overview}). It works over $\BFT$ and arbitrary trees and
transfers the result to the other tree classes and to $\BF$ at the end. We end
in the hole problem of the plain-entailment proof~\citep{Zhao26}, whose readers, holes and bounds
\cref{sec:bg} recalls: a \emph{reader} is a deterministic automaton with
rejected states and an admission relation on pairs of states; a
\emph{comparison} of a word $w$ is a pair of cuts $(s,e)$, $s<e$, such that the
suffix after $e$ is a prefix of the suffix after $s$; and a \emph{hole} is a
rejected word without an admitted comparison. That proof bounds a
shortest hole in three steps: a short rejected tail, a bound on the
\emph{depth}, a measure of nested repetition, and a compression of words of
bounded depth. What is new here is that some of our holes must also
\emph{contain} a comparison of a particular kind, and that a single word may
be dead because of a comparison with itself.

\begin{figure}[tbp]
\centering
\small
\hyphenpenalty=10000\relax
\begin{tikzpicture}[
    node distance=6mm,
    >=Stealth,
    stepbox/.style={draw, rounded corners=2pt, align=flush center,
                 text width=0.62\linewidth, inner sep=3pt, font=\small},
    ubox/.style={draw, rounded corners=2pt, align=flush center,
                 text width=0.25\linewidth, inner sep=3pt, font=\footnotesize},
    sidebox/.style={draw, dashed, rounded corners=2pt, align=flush left,
                 text width=0.43\linewidth, inner sep=3pt, font=\footnotesize},
    thm/.style={right, font=\footnotesize, inner sep=2pt}]
  \node[stepbox] (ent)
    {\textbf{Existential entailment}\\
     $\varphi\entails\Exists\psi$ over $\BFT$, any arity and variance, arbitrary
     trees, $\varphi$ and $\psi$ satisfiable};
  \node[stepbox, below=of ent] (ch)
    {\textbf{Leaf cones and the characterization}
     (\cref{sec:lf,sec:ch})\\
     \CondI\ the forced visible atoms of $\psi$ are forced by $\varphi$;
     \CondU\ no live seed derives a forbidden root of $\psi$;
     \CondII\ every visible $\psi$-consequence of a live seed lies in its
     $\varphi$-closure};
  \node[stepbox, below=of ch] (ex)
    {\textbf{Automata and holes} (\cref{sec:sm,sec:ex})\\
     \CondI\ is an inclusion of automata; under \CondU, \CondII\ says that no
     extracted two-word instance $I_{pq}$ has a hole, with self events and
     forced target heads};
  \node[stepbox, below=of ex] (nm)
    {\textbf{Three cases} (\cref{sec:nm})\\
     $I_{pq}$ has a hole iff it has a short hole, found through a branching
     state, a lasso or the pump, or a mortal hole};
  \node[stepbox, below=of nm] (mo)
    {\textbf{One separator word of bounded depth} (\cref{sec:mo})\\
     mortal holes $(u,t)$ are the $\Pi$-holes $u\sep t$ of one reader; their
     depth is bounded, and anchored compression bounds their length};
  \node[ubox, right=5mm of ex] (u)
    {\textbf{\CondU}: unit readers and the plain-entailment bounds
     (\cref{sec:un}, \cref{thm:un-decidable})};
  \draw[<->] (ent) -- node[thm]{\cref{thm:ch-characterization}} (ch);
  \draw[<->] (ch) -- node[thm]{\cref{thm:ex-reduction}} (ex);
  \draw[<->] (ex) -- node[thm]{\cref{thm:nm-three-cases}} (nm);
  \draw[<->] (nm) -- node[thm, align=left]{mortal branch,
    \cref{thm:mo-decide}} (mo);
  \draw[<->] (ch.east) -| (u.north);
  \node[sidebox, below=6mm of mo.south west, anchor=north west] (cls)
    {\textbf{Tree classes}: regular trees agree with arbitrary trees; finite
     trees add two finite-satisfiability tests
     (\cref{thm:tc-regular,thm:tc-finite}).};
  \node[sidebox, right=2mm of cls] (bf)
    {\textbf{The signature $\BF$}, $f$ covariant: a gadget on every variable
     reduces it to $\BFT$ (\cref{thm:tc-gadget}).};
\end{tikzpicture}
\caption{The route of the proof, for satisfiable $\varphi$ and $\psi$ over
$\BFT$ and arbitrary trees. Each double arrow is an equivalence, labeled with
the result that proves it, or for condition \CondU\ with its box; the last
vertical one concerns only the mortal branch of the case split, whose short
holes are found directly. Condition \CondU\ leaves the
main route: it is a hole problem of the kind solved for plain entailment,
decided by the bounds of \citet{Zhao26}. Here $p,q$ range over the visible signed variables. The
dashed boxes describe the transfers to the other tree classes and to $\BF$.
\Cref{sec:as} assembles the steps.}
\Description{Five boxes connected by double arrows, each labeled with a
theorem, lead from existential entailment to the characterization by three
conditions on leaf cones, to automata and two-word hole problems, to a case
split into short holes and mortal holes, and, for the mortal branch, to holes
of one separator word of bounded depth. A side box, reached from the
characterization, decides condition (U) by unit readers and the
plain-entailment bounds. Two dashed boxes describe the regular and finite tree classes and the
signature without top.}
\label{fig:intro-overview}
\end{figure}

\paragraph{Leaf cones}
Over $\BF$ with $f$ covariant, a tree is its set of constructor positions,
solution sets are closed under unions and intersections, and every literal
becomes a set of implications with one premise (\cref{rem:ch-one-sign}).
With $\Top$ and variance, a literal can relate a leaf $\Top$ of one tree to a
leaf $\Bot$ of another, and one set of positions no longer describes a tree.
We restore a Horn structure for every arity and variance by describing a tree
$t$ through
its two \emph{leaf cones}: the words whose path meets a leaf of $t$ of
\emph{sign}~$0$, and those whose path meets a leaf of sign~$1$, where the sign
of a leaf combines its label, $\Top$ or $\Bot$, with the polarity of its path.
The two cones are disjoint and closed under extension, every such pair of sets
is the pair of exactly one tree, and the order of variance $c$ adds leaves of
sign~$0$ and removes leaves of sign~$1$ (\cref{thm:lf-representation}). An
\emph{atom} $(v^\pm,w)$ says that $w$ lies in a cone of $v$. Every literal
becomes \emph{pure moves}, implications between atoms that copy, push or pop
the first letter of the word as on a stack, together with \emph{facts}, atoms
that every solution contains, and \emph{forbidden roots}, atoms that none
contains (\cref{lem:lf-clauses}). The only all-negative binary clauses are the
exclusions of opposite atoms, which are the same for every system. So least models are
closures of seeds, every atom of a closure has a single seed
(\cref{lem:lf-closure}), every pure move has a dual of opposite sign
(\cref{lem:lf-duality}), and two derivations that collide in an exclusion can
be merged into one (\cref{lem:lf-collision}).

\paragraph{Three conditions}
Collision return turns every obstruction to extending a solution of $\varphi$
into an obstruction at a single seed. If both systems are satisfiable,
$\varphi\entails\Exists\psi$ holds iff three conditions hold
(\cref{thm:ch-characterization}): \CondI, every visible atom that $\psi$
forces is forced by $\varphi$; \CondU, $\psi$ derives no forbidden root from a
visible atom that some solution of $\varphi$ contains; and \CondII, every
visible atom that $\psi$ derives from such an atom lies in the least solution
of $\varphi$ that contains it.

\paragraph{Pushdown summaries}
The pure moves form a pushdown system whose control states are the signed
variables. Two deterministic automata summarize it: after reading a word
$\gamma$, the \emph{upper automaton} of $p$ holds the signed variables reached
from $p$ by popping $\gamma$, and after reading $\beta$ the \emph{lower
automaton} of $q$ holds those from which $q$ is reached by pushing $\beta$.
Every pure derivation splits at its lowest point into a \emph{pop head}, read
by the first, and a \emph{push head}, read by the second, above a common
suffix (\cref{lem:sm-lowest-point}). Variance adds a new reason for an atom to
be dead. Besides a reachable forbidden root, or an opposite extension forced
by the facts, which make an atom \emph{dead by support}, the closure of a seed
$(p,w)$ may contain the opposite of an extension of the seed. For a seed that is not already dead
by support, this happens exactly when two cuts of the single word $w$ form a
comparison whose automaton states are related: a \emph{self event} (\cref{thm:sm-dead}). Self events occur only
through contravariant directions (\cref{rem:sm-parity}), and the set of live
words need not be regular (\cref{ex:sm-arrow}). Condition \CondI\ is an
inclusion of automata, decidable in polynomial space (\cref{prop:sm-decide}).

\paragraph{Condition (U) as a hole problem}
A failure of \CondU\ is a word $w$ such that $(p,w)$ is live in $\varphi$ and
leads to a forbidden root in $\psi$. Liveness is the absence of deadness by
support, a regular property, and of self events, which are comparisons of $w$.
So the failures are exactly the holes of a single reader, the \emph{unit
reader}, whose admitted comparisons are the self events
(\cref{lem:un-holes}). It has the rejected-path property of \citet{Zhao26}, and
the bounds proved there decide \CondU\ (\cref{thm:un-decidable}). A smaller reader
with the same holes, together with a clock that reads only the antecedent and
bounds the depth, keeps the length bound for \CondU\ below the one for mortal
holes.

\paragraph{Two-word holes with self events}
A failure of \CondII\ consists of a live seed $(p,t)$, a pure derivation of
$\psi$ from it to a visible atom $(q,u)$, and the fact that $(q,u)$ lies
outside the closure of the seed in $\varphi$. Each part is read by automata on
the two words: the derivation is a \emph{head event}, a pair of cuts at which
$t$ and $u$ have equal suffixes; membership in the closure is an \emph{event},
a pair of cuts $(s,e)$ at which the suffix of $t$ after $e$ is a prefix of the
suffix of $u$ after $s$; and liveness of the seed forbids self events of $t$.
So, once \CondU\ holds, \CondII\ fails iff an \emph{extracted instance}
$I_{pq}$ of a hole problem for two words has a \emph{hole}: a safe spine $u$,
a rejected target $t$, a head event, and neither an event nor a self event
(\cref{def:ex-ss2self,thm:ex-reduction}). The target automaton marks a cut of
the seed as rejected only while the seed is still a candidate, neither dead by
support in $\varphi$ nor root-dead in $\psi$. Under \CondU\ this enriched
rejection loses nothing, and it forces every pop head of $\psi$ to end before
the first rejected cut, condition (FH) (\cref{lem:ex-fh}).

\paragraph{Holes that are not mortal}
Cut a hole $(u,t)$ at the first rejected cut of its target. By (FH) the rest
$\theta$ of the target is a common suffix of both words, and what remains is a
\emph{pre-hole}. Appending $\theta$ to a pre-hole creates no new event and no
new self event, and it keeps an old one exactly when $\theta$ continues its
\emph{overhang}, the part of one suffix that the other does not cover,
periodically (\cref{lem:nm-continuation}). So the question is whether one safe
continuation breaks every such periodicity obligation. If the spine ends
in a \emph{branching} state, which reaches a cycle with a choice of letters, a
breaking continuation exists by the theorem of \citet{FineWilf65}
(\cref{lem:nm-breaker}). If it ends in a \emph{thin} state, long safe
continuations follow finitely many lassos; a lasso that is not purely periodic
breaks everything, and for a purely periodic lasso the \emph{pump} shortens
the pre-hole itself, keeping the obligations broken (\cref{thm:nm-pump}).
Short pre-holes are found by a finite graph (\cref{prop:nm-short-prehole}). In
the remaining case the hole is \emph{mortal}: its spine ends in a state
without an infinite safe continuation, and its target ends shortly after its
first rejected cut (\cref{thm:nm-three-cases}). Outside the mortal case,
holes are polynomial in the sizes of the automata, hence singly exponential
in~$s$.

\paragraph{Mortal holes}
In a mortal hole the continuation is short, but the common suffix of the two
words can be long, so the pair itself must be shortened, with both words
together. We pack them into one word $u\sep t$ around a separator. Events and
self events become comparisons of this word, and a head event becomes an
\emph{anchored} comparison, one whose matched factor is followed by the
separator (\cref{lem:mo-sep}). Mortal holes are exactly the
\emph{$\Pi$-holes} of one reader, its holes with an anchored comparison whose
state pair lies in a relation $\Pi$ (\cref{lem:mo-pi-reader}), and the short
rejected tail bounds their depth by the depth theorem of \citet{Zhao26}
(\cref{lem:mo-depth}). The bounded-depth compression of that paper keeps every negative
condition but can delete the letters on which the anchored comparison depends
(\cref{ex:mo-loses}). A \emph{positive channel} records the anchored witness
at each level of the hierarchy (\cref{lem:mo-desub}). The channel is added to
the information that this compression compares before it merges
letters or deletes columns, so every shortening keeps some anchored witness.
This gives the anchored compression theorem (\cref{thm:mo-anchored}) and a
triply exponential bound on mortal holes (\cref{cor:mo-extracted}), the
largest of our bounds, which decides the two-word problem under (FH)
(\cref{thm:mo-decide}).

\paragraph{Tree classes and the signature without top}
Over regular trees the closure of regular seeds is regular, and a failure has
at most two seeds, so regular and arbitrary trees give the same answer
(\cref{thm:tc-regular}). Least models are infinite in general, so finite trees
need a different argument: the at most two observations behind a failure are
expressed by finite ground bounds, and adding finite ground bounds preserves
finite satisfiability, for every variance (\cref{thm:tc-pinning}). This gives
the dispatch of \cref{thm:tc-finite}. For $\BF$ with $f$ covariant, the
literals $v\le y_v$ and $y_v\le f(y_v,\ldots,y_v)$ force $v$ to be free of
$\Top$, and placing them on every visible and every hidden variable reduces
$\BF$ to $\BFT$ in each tree class (\cref{thm:tc-gadget}).

\paragraph{The running examples}
Two examples recur through the paper (\cref{tab:intro-examples}). The first is
the shift closure of \cref{ex:intro-shift}, for a covariant constructor; the
second follows the antecedent $x\le x\to x$ for the arrow, in which $x$ occurs
at both polarities. For a covariant constructor, and more generally on
instances with a joint parity assignment, self events do not occur and every
dead atom is dead by support
(\cref{rem:ch-parity,rem:sm-parity,rem:ex-covariant}); over $\BF$ the three
conditions compare unions of solutions, least solutions, and closures of
single positions (\cref{rem:ch-one-sign}).

\begin{table}[tbp]
\caption{The running examples.}
\label{tab:intro-examples}
\centering
\small
\begin{tabular}{@{}>{\raggedright\arraybackslash}p{0.21\linewidth}>{\raggedright\arraybackslash}p{0.17\linewidth}>{\raggedright\arraybackslash}p{0.54\linewidth}@{}}
\toprule
Example & Sections & What it shows\\
\midrule
\multicolumn{3}{@{}l}{\emph{Shift closure}: $\psi=x\le z\wedge z\le y\wedge
  z\le f(z,z)$, with $f$ binary and covariant}\\
\cref{ex:intro-shift} & \cref{sec:intro} & an existential entailment that is
  an infinite family of plain entailments\\
\cref{ex:sm-running} & \cref{sec:sm} & the signed system over $\BFT$ and its
  pure derivations\\
\cref{ex:un-running} & \cref{sec:un} & \CondU\ fails at $x=y=\Top$, and holds
  once $y$ is restricted to trees without $\Top$\\
\cref{ex:ex-running} & \cref{sec:ex} & \CondII\ fails by the hole $(01,1)$ of
  an extracted instance, whose spine ends in a branching state\\
\cref{ex:nm-running} & \cref{sec:nm} & the pre-hole $(0,\eps)$, its single
  obligation, and, after \cref{lem:nm-breaker}, a breaker\\
\cref{ex:as-running} & \cref{sec:as} & the procedure end to end: the hole word
  $01\sep1$, its two pins, and the finite countermodel of
  \cref{ex:intro-shift}\\
\midrule
\multicolumn{3}{@{}l}{\emph{The arrow}: the antecedent $x\le x\to x$}\\
\cref{ex:lf-clauses,ex:lf-least} & \cref{sec:lf} & a single-seed
  inconsistency\\
\cref{ex:ch-arrow} & \cref{sec:ch} & the entailment
  $x\le x\to x\entails\exists z\,(x\le z\to z)$, whose failure without the
  antecedent needs collision return\\
\cref{ex:sm-arrow} & \cref{sec:sm} & a nonregular set of dead words\\
\cref{ex:un-arrow,ex:ex-arrow} & \cref{sec:un,sec:ex} & why admitted self
  events and the enrichment of the target automaton are needed\\
\bottomrule
\end{tabular}
\end{table}

\paragraph{Contributions}
\begin{itemize}[leftmargin=*]
\item \emph{Decidability} (\cref{thm:intro-main}). Existential entailment is
  decidable for one constructor of every arity and every variance over
  $\BFT$, and over $\BF$ with $f$ covariant, in all three tree classes. This
  settles the decidability question of \citet{TrifonovSmith96} for semantic
  subsumption of constrained types.
\item \emph{Signed leaf cones} (\cref{sec:lf,sec:ch}). A representation of
  trees of any variance in which every system is a Horn system with unary
  rules, and a characterization of existential entailment by three conditions
  on single seeds.
\item \emph{Self events} (\cref{sec:sm,sec:un,sec:ex}). Deadness is deadness by
  support or a self event, a comparison of one word with itself. Both
  \CondU\ and \CondII\ become hole problems whose admissions include self
  events, the second with an enriched rejection that makes target heads forced.
\item \emph{Non-mortal holes} (\cref{sec:nm}). Breakers, lassos and a pump
  with polynomial capacity give singly exponential witnesses outside the
  mortal case.
\item \emph{Anchored compression} (\cref{sec:mo}). Positive channels extend the
  bounded-depth compression of \citet{Zhao26} to holes that must contain
  an anchored comparison.
\item \emph{Tree classes and $\BF$} (\cref{sec:tc}). Regular and finite
  transfers for every variance, and a reduction of covariant $\BF$ to $\BFT$.
\item \emph{Complexity} (\cref{sec:cx}). Elementary witness bounds; a
  2-EXPSPACE decision by summary search; PSPACE for $n\le1$; exponential
  shortest witnesses; and PSPACE-hardness for every fixed arity $n\ge2$,
  already with $f$ covariant.
\item \emph{Applications} (\cref{sec:ap}). Subsumption of constrained type
  schemes and correctness of simplification are decidable, including arrow
  types without any restriction on polarity, and several classes reduce
  directly to plain entailment.
\item \emph{Mechanization} (\cref{sec:mechanization}). Computable decision
  procedures for all six cases of \cref{thm:intro-main}, checked in Lean~4,
  including anchored compression with an elementary bound.
\end{itemize}

\paragraph{Organization}
\Cref{sec:bg} fixes notation and recalls the results of \citet{Zhao26} that we use.
\Cref{sec:lf,sec:ch,sec:sm,sec:un,sec:ex,sec:nm,sec:mo,sec:tc} follow the
route above, and \cref{sec:as} assembles the proof.
\Cref{sec:ap,sec:cx,sec:mechanization,sec:rel,sec:concl} discuss
applications, complexity, the mechanization, related work, and open problems. Omitted proofs are in the appendices.
\Cref{sec:lf,sec:ch} give the characterization, and
\cref{sec:sm,sec:un,sec:ex} turn it into hole problems; \cref{sec:nm,sec:mo} are the technical core, and
\cref{sec:mo} assumes \cref{sec:bg-hierarchy}. A reader interested in the
result and its use may go from \cref{sec:ch} to \cref{sec:as,sec:ap}.
